% Archive record 23, line 1.
\hypertarget{part-2-23-0001-ii-e-data-leakage-and-distribution-shift}{%
\section{II-E --- Data Leakage and Distribution Shift}\label{part-2-23-0001-ii-e-data-leakage-and-distribution-shift}}

此前 II-A 至 II-D，我们主要研究的是：在数据生成分布相对稳定的条件下，为什么模型选择与自适应反馈会造成过拟合。

现在开始放松这一假设。

对于 X Program \#161，观察到 OOS Performance Decay 之后，至少存在三种不同的解释：

\begin{readertable}
\begin{longtable}[]{@{}
  >{\raggedright\arraybackslash}p{(\columnwidth - 2\tabcolsep) * \real{0.33}}
  >{\raggedright\arraybackslash}p{(\columnwidth - 2\tabcolsep) * \real{0.63}}@{}}
\toprule
Mechanism & 核心问题 \\
\midrule
\endhead
\textbf{Adaptive Overfitting} & Agent 过度适应了有限的历史评估数据 \\
\textbf{Data Leakage} & Agent 使用了正常预测或评估过程中不应获得的信息 \\
\textbf{Distribution Shift} & Agent 学到的规律是真实的，但目标环境已经发生变化 \\
\bottomrule
\end{longtable}
\end{readertable}

这三种机制可能同时发生，却需要不同的诊断方法和解决方案。

本阶段的重点不是增加大量数学证明，而是建立一套能够区分这些机制的 \textbf{Diagnostic Framework}。尤其需要理解：为什么一条 IS 上升、OOS 下降的曲线，本身不足以说明性能退化的原因。

\hypertarget{part-2-23-0001-data-leakageux6570ux636eux6cc4ux6f0f}{%
\subsection{1. Data Leakage（数据泄漏）}\label{part-2-23-0001-data-leakageux6570ux636eux6cc4ux6f0f}}

\hypertarget{part-2-23-0001-definition}{%
\subsubsection{1.1 Definition}\label{part-2-23-0001-definition}}

Data Leakage 指模型开发、训练或者评估过程使用了在相应的真实应用场景中不应当获得的信息，从而破坏评估的有效性。

这里最重要的不是信息是否能够被获取，而是：

\textbf{在相应的预测时刻和实验协议下，使用该信息是否合法。}

例如，我们希望在时刻 \(t\) 预测下一交易日的收益率：
\[r_{t+1}=g(X_t)+\varepsilon_{t+1}.\]

设 \(\mathcal F_t\) 表示截至时刻 \(t\) 合法获得的信息所生成的 \(\sigma\)-algebra。

那么预测结果应当满足：
\[\boxed{
\hat r_{t+1}=h(X_t)
\quad\text{is }\mathcal F_t\text{-measurable}.
}\]

这意味着预测必须基于当前已经可以获得的信息。

如果 \(X_t\) 中包含只有在 \(t+1\) 时刻才能获得的变量，就违反了这一条件。

这正是金融领域 \textbf{Look-Ahead Bias（前视偏差）} 的数学本质。

\hypertarget{part-2-23-0001-data-leakage-ux7684ux4e3bux8981ux7c7bux578b}{%
\subsubsection{1.2 Data Leakage 的主要类型}\label{part-2-23-0001-data-leakage-ux7684ux4e3bux8981ux7c7bux578b}}

\begin{readertable}
\begin{longtable}[]{@{}
  >{\raggedright\arraybackslash}p{(\columnwidth - 2\tabcolsep) * \real{0.35}}
  >{\raggedright\arraybackslash}p{(\columnwidth - 2\tabcolsep) * \real{0.61}}@{}}
\toprule
Leakage Type & 具体例子 \\
\midrule
\endhead
Target Leakage & 使用未来收益率构建当前预测特征 \\
Temporal Leakage & 使用预测时刻之后发布的财务报表 \\
Preprocessing Leakage & 在划分训练集和测试集之前，使用全量数据进行标准化 \\
Evaluation Leakage & 将最终测试集得分反馈给 Agent \\
Cross-sample Leakage & 训练集和测试集包含重复或高度重叠的样本 \\
Benchmark Contamination & LLM 在预训练阶段已经接触过评估任务及其答案 \\
\bottomrule
\end{longtable}
\end{readertable}

这里有一个值得注意的细节：Data Leakage 并不一定导致模型参数发生变化。

只要它改变了模型的选择、特征构造、Agent Revision 或最终实验结论，就可能损害评估的独立性。

\hypertarget{part-2-23-0001-ux4e00ux4e2aux5177ux4f53ux7684ux91d1ux878dux4f8bux5b50}{%
\subsubsection{1.3 一个具体的金融例子}\label{part-2-23-0001-ux4e00ux4e2aux5177ux4f53ux7684ux91d1ux878dux4f8bux5b50}}

假设我们构造五日移动平均线（Moving Average）作为金融预测特征。

\textbf{Valid Feature --- Trailing Moving Average}

\bookgraphic{0.830}{part-2-23-0001-001-01}{}

\[\mathrm{MA}_t
  =
  \frac15\sum_{j=0}^{4}P_{t-j}\]

假设在时刻 t 的收盘价格已知之后进行预测。

\textbf{Invalid Feature --- Centered Moving Average}

\bookgraphic{0.830}{part-2-23-0001-001-02}{}

\[\mathrm{MA}_t^{\mathrm{centered}}
  =
  \frac15\sum_{j=-2}^{2}P_{t+j}\]

包含未来两个交易日的价格，不能直接作为时刻 t 的实时预测特征。

在 Python 中，两者可能只相差一个参数：

\begin{verbatim}
# Valid, assuming today's close is available
df["ma"] = df["close"].rolling(5).mean()

# Invalid for real-time prediction at t
df["ma"] = df["close"].rolling(5, center=True).mean()
\end{verbatim}

对于一个可以自主修改 Python 代码的 RSI Agent，后一种修改可能显著改善历史预测得分，进而被选择机制保留下来。

这说明 \textbf{Revision Classification 不仅需要检查模型结构是否更加复杂，还需要检查修改是否改变了信息的可用性。}

\hypertarget{part-2-23-0001-data-leakage-ux4e0e-adaptive-overfitting-ux7684ux533aux522b}{%
\subsection{2. Data Leakage 与 Adaptive Overfitting 的区别}\label{part-2-23-0001-data-leakage-ux4e0e-adaptive-overfitting-ux7684ux533aux522b}}

这两个概念存在交集，但不完全等价。

\begin{readertable}
\begin{longtable}[]{@{}
  >{\raggedright\arraybackslash}p{(\columnwidth - 4\tabcolsep) * \real{0.29}}
  >{\raggedright\arraybackslash}p{(\columnwidth - 4\tabcolsep) * \real{0.31}}
  >{\raggedright\arraybackslash}p{(\columnwidth - 4\tabcolsep) * \real{0.37}}@{}}
\toprule
Dimension & Data Leakage & Adaptive Overfitting \\
\midrule
\endhead
根本机制 & 使用不应获得的信息 & 过度适应有限评估数据 \\
是否需要反复搜索 & 不需要 & 通常涉及数据驱动的自适应过程 \\
单次实验能否发生 & 可以 & 典型情形涉及多轮反馈 \\
是否必然违反实验协议 & 是，按这里采用的定义 & 不一定 \\
主要解决方式 & 修正信息访问与评估协议 & 控制搜索复杂度及反馈信息量 \\
\bottomrule
\end{longtable}
\end{readertable}

例如，Agent 可以完全遵守开发集的访问规则，却因为反复选择而对开发集过拟合。

这种情况下存在 Adaptive Overfitting，但不一定存在 Data Leakage。

反过来，一个完全没有经过迭代优化的模型，如果直接使用了未来收益率作为特征，也会产生 Data Leakage。

而如果 Agent 反复获得本应封闭的最终测试集得分，那么两种机制可能同时存在。

\textbf{Research Implication：} 对于 X Program \#161，我们应当首先通过代码审计排除明显的 Data Leakage，否则后续对 Adaptive Overfitting 的统计分析可能建立在错误的实验基础上。

\hypertarget{part-2-23-0001-distribution-shiftux5206ux5e03ux504fux79fb}{%
\subsection{3. Distribution Shift（分布偏移）}\label{part-2-23-0001-distribution-shiftux5206ux5e03ux504fux79fb}}

\hypertarget{part-2-23-0001-definition-1}{%
\subsubsection{3.1 Definition}\label{part-2-23-0001-definition-1}}

沿用此前的记号，设开发数据来自总体分布 \(P\)，而目标环境中的数据来自总体分布 \(Q\)。

经典 Statistical Learning Theory 通常假设：
\[P(X,Y)=Q(X,Y).\]

但在实际应用中，这一假设经常不成立：
\[\boxed{P(X,Y)\neq Q(X,Y)}\]

这就是 Distribution Shift，亦称 Dataset Shift。

由于联合分布可以进行如下分解：
\[P(X,Y)=P(Y\mid X)P(X)\]

以及：
\[P(X,Y)=P(X\mid Y)P(Y),\]

我们可以进一步研究究竟是哪一个组成部分发生了变化。

一种经典分类方式是将 Distribution Shift 分为 Covariate Shift、Prior Probability Shift 和 Concept Shift。

\href{https://www.sciencedirect.com/science/article/pii/S0031320311002901}{来源 1}

\hypertarget{part-2-23-0001-covariate-shiftux534fux53d8ux91cfux504fux79fb}{%
\subsubsection{3.2 Covariate Shift（协变量偏移）}\label{part-2-23-0001-covariate-shiftux534fux53d8ux91cfux504fux79fb}}

定义：
\[\boxed{
\begin{aligned}
P(X)&\neq Q(X),\\
P(Y\mid X)&=Q(Y\mid X).
\end{aligned}
}\]

即输入变量的分布发生变化，但输入与输出之间的条件分布保持不变。

\bookgraphic{0.558}{part-2-23-0001-003-01}{}

\[P(Y\mid X)=Q(Y\mid X)\]

输入数据的分布发生变化，但给定输入后的预测规律保持不变。

\textbf{Financial Example}

假设 \(X\) 表示市场波动率，\(Y\) 表示某个资产的未来收益率。

开发数据主要来自低波动率时期，而未来数据主要来自高波动率时期。但是，在给定市场波动率后，收益率的条件分布没有改变。

此时就属于 Covariate Shift。

尽管预测规律没有变化，模型仍然可能表现变差，因为它在训练阶段很少接触高波动率数据。

\hypertarget{part-2-23-0001-ux6570ux5b66ux5904ux7406importance-weighting}{%
\paragraph{数学处理：Importance Weighting}\label{part-2-23-0001-ux6570ux5b66ux5904ux7406importance-weighting}}

定义两个分布下的 Population Risk：
\[R_P(h)=\mathbb E_P[\ell(h(X),Y)],\]

\[R_Q(h)=\mathbb E_Q[\ell(h(X),Y)].\]

在 Covariate Shift、目标输入分布对源输入分布绝对连续等条件下：
\[\begin{aligned}
R_Q(h)
&=\int \ell(h(x),y)q(x,y)\,\mathrm dx\,\mathrm dy\\
&=\int \ell(h(x),y)
\frac{q(x)}{p(x)}
p(x,y)\,\mathrm dx\,\mathrm dy.
\end{aligned}\]

因此：
\[\boxed{
R_Q(h)
=
\mathbb E_P
\left[
\frac{q(X)}{p(X)}
\ell(h(X),Y)
\right]
}\]

其中：
\[w(X)=\frac{q(X)}{p(X)}\]

称为 Importance Weight。

其基本思想是：对未来环境中更常出现的输入赋予更高权重。

但如果两个分布的重叠程度很低，Importance Weight 可能非常大，进而增加估计方差。如果未来出现了源分布完全不覆盖的输入，这个方法也无法直接解决问题。

\hypertarget{part-2-23-0001-label-shiftux6807ux7b7eux504fux79fb}{%
\subsubsection{3.3 Label Shift（标签偏移）}\label{part-2-23-0001-label-shiftux6807ux7b7eux504fux79fb}}

定义：
\[\boxed{
\begin{aligned}
P(Y)&\neq Q(Y),\\
P(X\mid Y)&=Q(X\mid Y).
\end{aligned}
}\]

即标签的边际分布发生变化，但给定标签后的特征分布保持不变。

假设 Y 为市场上涨或下跌的二分类标签。

Source \(P\)

70\%

30\%

Target \(Q\)

35\%

65\%

Market Up

Market Down

\[P(X\mid Y)=Q(X\mid Y)\]

例如，假设在开发环境中，70\% 的交易日上涨；而在未来环境中，只有 35\% 的交易日上涨。

如果给定市场涨跌状态后的特征分布保持不变，就构成 Label Shift。

这里可以直接联系 Bayes' Theorem：
\[P(Y\mid X)
=
\frac{P(X\mid Y)P(Y)}{P(X)}.\]

即使 \(P(X\mid Y)\) 保持不变，改变 Prior Probability \(P(Y)\) 也可能改变 Posterior Probability \(P(Y\mid X)\)。

因此，Label Shift 可能导致原有分类概率与分类阈值不再适合目标环境。

2018 年 Lipton、Wang 和 Smola 在 ICML 发表的研究讨论了如何在没有目标环境标签的情况下，估计并校正满足特定条件的 Label Shift。

\href{https://proceedings.mlr.press/v80/lipton18a}{来源 1}

在本课题中，Label Shift 的优先级相对较低，尤其是在研究对象为连续收益率预测而非分类任务时。

\hypertarget{part-2-23-0001-concept-driftux6982ux5ff5ux6f02ux79fb}{%
\subsubsection{3.4 Concept Drift（概念漂移）}\label{part-2-23-0001-concept-driftux6982ux5ff5ux6f02ux79fb}}

这是金融场景中尤其值得关注的一种 Distribution Shift。

其核心特征为：
\[\boxed{
P(Y\mid X)\neq Q(Y\mid X)
}\]

即输入与输出之间的统计关系发生变化。

为了单独研究这种变化，可以考虑保持输入分布不变的 Pure Concept Shift：
\[P(X)=Q(X),
\qquad
P(Y\mid X)\neq Q(Y\mid X).\]

Concept Drift 通常强调这种关系随时间发生变化。不同文献的术语定义略有差异，我们在本课程中采用上述条件分布变化的含义。

\bookgraphic{0.484}{part-2-23-0001-006-01}{合成数据示意：输入变量的分布没有明显变化，但其与预测目标之间的关系发生反转。}

例如，Momentum Strategy 在趋势性市场中可能有效，但在均值回归明显的市场中可能失效。

与 Covariate Shift 不同，这里的困难不是某些输入出现得更加频繁，而是原有的预测规则本身可能不再成立。

因此，仅仅通过对旧数据进行 Importance Weighting，一般不足以解决 Concept Drift。

\hypertarget{part-2-23-0001-ux4e09ux79cd-distribution-shift-ux7684ux5bf9ux6bd4}{%
\subsubsection{3.5 三种 Distribution Shift 的对比}\label{part-2-23-0001-ux4e09ux79cd-distribution-shift-ux7684ux5bf9ux6bd4}}

\begin{readertable}
\begin{longtable}[]{@{}
  >{\raggedright\arraybackslash}p{(\columnwidth - 6\tabcolsep) * \real{0.25}}
  >{\raggedright\arraybackslash}p{(\columnwidth - 6\tabcolsep) * \real{0.26}}
  >{\raggedright\arraybackslash}p{(\columnwidth - 6\tabcolsep) * \real{0.22}}
  >{\raggedright\arraybackslash}p{(\columnwidth - 6\tabcolsep) * \real{0.23}}@{}}
\toprule
& Covariate Shift & Label Shift & Pure Concept Shift \\
\midrule
\endhead
发生变化 & \(P(X)\) & \(P(Y)\) & \(P(Y\mid X)\) \\
保持不变 & \(P(Y\mid X)\) & \(P(X\mid Y)\) & \(P(X)\) \\
核心含义 & 输入环境变化 & 类别比例变化 & 预测规律变化 \\
候选处理方式 & Importance Weighting & Prior Correction & Model Updating \\
金融任务学习优先级 & P1 & P2 & P0 \\
\bottomrule
\end{longtable}
\end{readertable}

这些是具有特定假设的理想化分类，不是互斥且穷尽所有情况的划分。真实金融环境中的分布变化往往是多种机制共同作用的结果。

\hypertarget{part-2-23-0001-mathematical-example-ux6ca1ux6709ux8fc7ux62dfux5408ux4e3aux4ec0ux4e48ux4ecdux7136ux4f1aux9000ux5316}{%
\subsection{4. Mathematical Example --- 没有过拟合，为什么仍然会退化？}\label{part-2-23-0001-mathematical-example-ux6ca1ux6709ux8fc7ux62dfux5408ux4e3aux4ec0ux4e48ux4ecdux7136ux4f1aux9000ux5316}}

我们利用 Linear Regression 构造一个非常简单的反例。

假设开发环境中的真实数据生成机制为：
\[Y=\beta X+\varepsilon,
\qquad
X\sim\mathcal N(0,1),\]

其中 \(\varepsilon\) 与 \(X\) 独立，且：
\[\varepsilon\sim\mathcal N(0,\sigma^2).\]

现在假设未来环境的真实规律变为：
\[Y=(\beta+\Delta)X+\varepsilon.\]

其中 \(\Delta\) 表示预测关系的变化程度。

为了完全排除有限样本过拟合，我们假设 Agent 已经准确知道开发环境中的真实参数 \(\beta\)。 因此：
\[h(X)=\beta X.\]

在开发环境中的 Population Risk 为：
\[R_P(h)=\sigma^2.\]

但在目标环境中：
\[\begin{aligned}
R_Q(h)
&=\mathbb E[(Y-h(X))^2]\\
&=\mathbb E[(\Delta X+\varepsilon)^2]\\
&=\sigma^2+\Delta^2.
\end{aligned}\]

因此：
\[\boxed{
R_Q(h)-R_P(h)=\Delta^2
}\]

\textbf{Concept Drift --- Risk Simulation}

假设 \(\beta=1,\ \sigma^2=0.25\)，模型准确掌握开发环境中的预测规律。

Concept Shift \(\Delta\)\\
\textbf{0.50}

0.5（范围 -1.5--1.5）

\bookgraphic{0.573}{part-2-23-0001-007-01}{}

Learned model

Target relationship

Source Risk

\textbf{0.250}

Target Risk

\textbf{0.500}

Shift-induced Loss

\textbf{+0.250}

shift = -1.5

\bookgraphic{0.573}{part-2-23-0001-007-state-01}{}

\textbf{数学性质：} 在上述假设下严格成立。

这个例子的意义在于，即使模型完全没有发生 Statistical Overfitting，目标环境中的性能仍然可以显著下降。

进一步，如果 Agent 在递归改进过程中越来越精确地学习源环境中特有的规律，那么它的源环境表现可能持续提高，而目标环境表现反而下降。

但这需要分析不同版本对源环境和目标环境的适应程度，不能仅由上述单模型例子直接推出。

\hypertarget{part-2-23-0001-decomposing-the-generalization-gap}{%
\subsection{5. Decomposing the Generalization Gap}\label{part-2-23-0001-decomposing-the-generalization-gap}}

我们在 Part I 中已经介绍过基本分解，现在进一步完善它。

\hypertarget{part-2-23-0001-population-level-decomposition}{%
\subsubsection{5.1 Population-level Decomposition}\label{part-2-23-0001-population-level-decomposition}}

设 \(D_n\) 来自源分布 \(P\)。

沿用 Risk Notation：
\[\boxed{
\begin{aligned}
R_Q(h)-\widehat R_n(h)
={}&
\underbrace{
R_P(h)-\widehat R_n(h)
}_{\text{Statistical Generalization Gap}}\\
&+
\underbrace{
R_Q(h)-R_P(h)
}_{\text{Distribution Shift Gap}}.
\end{aligned}
}\]

定义：
\[G_{\mathrm{stat}}(h)
=
R_P(h)-\widehat R_n(h),\]

\[G_{\mathrm{shift}}(h)
=
R_Q(h)-R_P(h).\]

那么：
\[\boxed{
G_{\mathrm{total}}(h)
=
G_{\mathrm{stat}}(h)+G_{\mathrm{shift}}(h)
}\]

这一恒等式将有限样本过拟合与环境变化造成的性能差距分离。

\hypertarget{part-2-23-0001-empirical-evaluation-ux4e2dux7684ux7b2cux4e09ux9879}{%
\subsubsection{5.2 Empirical Evaluation 中的第三项}\label{part-2-23-0001-empirical-evaluation-ux4e2dux7684ux7b2cux4e09ux9879}}

但是真实科研中，我们通常也无法直接计算 \(R_Q(h)\)，只能使用独立的 OOS Dataset 进行估计。

设该数据集为 \(D_m^Q\)，对应经验风险：
\[\widehat R_m^Q(h).\]

于是，真正能够观察到的 IS/OOS Gap 为：
\[\boxed{
\begin{aligned}
\widehat R_m^Q(h)-\widehat R_n(h)
={}&
\underbrace{
R_P(h)-\widehat R_n(h)
}_{\text{Statistical Gap}}\\
&+
\underbrace{
R_Q(h)-R_P(h)
}_{\text{Distribution Shift}}\\
&+
\underbrace{
\widehat R_m^Q(h)-R_Q(h)
}_{\text{OOS Evaluation Noise}}.
\end{aligned}
}\]

\textbf{Observed Gap Decomposition}

假想的风险数据：越低越好。

\bookgraphic{0.486}{part-2-23-0001-008-01}{实际观察到的 Gap 为 0.23，但并不能将其全部归因于过拟合。}

这个分解比前面的两项分解更加接近真实科研环境。

但它仍然只是一个代数恒等式。由于 Population Risk 通常不可观测，我们无法仅凭两个经验得分准确估计其中每一项。

另外，Data Leakage 没有在这里作为独立的第四项出现，因为 Leakage 可能直接改变模型、数据或评估过程，使这些风险量不再具有原本预期的解释。它需要通过实验协议和代码审计另行处理。

\hypertarget{part-2-23-0001-distribution-shift-ux4e0e-rsi-performance-decay}{%
\subsection{6. Distribution Shift 与 RSI Performance Decay}\label{part-2-23-0001-distribution-shift-ux4e0e-rsi-performance-decay}}

这里需要区分两个不同的研究维度：

\begin{itemize}
\tightlist
\item
  \textbf{Across Iterations}：随着 Agent 版本变化，模型性能是否退化？
\item
  \textbf{Across Environments}：随着评估环境变化，模型性能是否退化？
\end{itemize}

我们假设有三个版本，并在两个固定环境中分别评价它们。

\textbf{Cross-environment Performance Evaluation}

假想的 Performance Score，越高越好。

\bookgraphic{0.723}{part-2-23-0001-009-01}{}

在这个例子中，从 \(A_1\) 到 \(A_2\)：

\begin{readertable}
\begin{longtable}[]{@{}
  >{\raggedright\arraybackslash}p{(\columnwidth - 2\tabcolsep) * \real{0.57}}
  >{\raggedleft\arraybackslash}p{(\columnwidth - 2\tabcolsep) * \real{0.39}}@{}}
\toprule
Evaluation Environment & Performance Change \\
\midrule
\endhead
IS & +0.20 \\
Independent Same-distribution & −0.03 \\
Shifted Environment & −0.25 \\
\bottomrule
\end{longtable}
\end{readertable}

这就说明，Agent 的后期修改可能同时存在两个问题：

\begin{enumerate}
\def\labelenumi{\arabic{enumi}.}
\tightlist
\item
  即使在同分布的独立数据上，泛化性能也出现了一定下降。
\item
  在发生分布变化的环境中，这种下降更加严重。
\end{enumerate}

对于固定目标分布 \(Q\)，所有版本都面对同一个目标环境。因此，如果不同版本的 OOS 表现发生变化，说明它们对该环境的适应能力存在差异，而不是说金融市场随着 RSI 的每轮迭代必然发生了变化。

这一区分对于设计实验非常重要。

\hypertarget{part-2-23-0001-experimental-design-ux5982ux4f55ux533aux5206ux4e09ux79cdux9000ux5316ux673aux5236}{%
\subsection{7. Experimental Design --- 如何区分三种退化机制？}\label{part-2-23-0001-experimental-design-ux5982ux4f55ux533aux5206ux4e09ux79cdux9000ux5316ux673aux5236}}

这是本阶段与 X Program \#161 联系最直接的部分。

我们可以构建三个评估环境：

\textbf{RSI Agent Versions}

\begin{verbatim}
$A_0,A_1,\ldots,A_K$
\end{verbatim}

↓

\textbf{IS}

\begin{verbatim}
    $D_{\mathrm{IS}}$
\end{verbatim}

反馈给 Agent

\textbf{OOS-P}

\begin{verbatim}
    $D_{\mathrm{OOS}}^P$
\end{verbatim}

独立、同分布

\textbf{OOS-Q}

\begin{verbatim}
    $D_{\mathrm{OOS}}^Q$
\end{verbatim}

独立、不同分布

两个 OOS 环境均不参与 Agent 的版本生成和选择。 其中：

\begin{itemize}
\tightlist
\item
  \(D_{\mathrm{IS}}\) 来自分布 \(P\)，负责驱动 RSI。
\item
  \(D_{\mathrm{OOS}}^P\) 也来自 \(P\)，但独立于优化过程。
\item
  \(D_{\mathrm{OOS}}^Q\) 来自不同分布 \(Q\)，负责检验跨环境泛化能力。
\end{itemize}

\hypertarget{part-2-23-0001-ux56dbux79cdux6f5cux5728ux5b9eux9a8cux7ed3ux679c}{%
\subsubsection{7.1 四种潜在实验结果}\label{part-2-23-0001-ux56dbux79cdux6f5cux5728ux5b9eux9a8cux7ed3ux679c}}

\textbf{Case A --- Genuine Improvement}

\bookgraphic{0.453}{part-2-23-0001-011-01}{三个环境均改善，存在真实泛化收益的证据。}

\textbf{Case B --- Overfitting}

\bookgraphic{0.453}{part-2-23-0001-011-02}{同分布样本外表现也下降，应重点检查统计过拟合。}

\textbf{Case C --- Shift Sensitivity}

\bookgraphic{0.453}{part-2-23-0001-011-03}{同分布泛化继续改善，但跨环境泛化能力下降。}

\textbf{Case D --- Potential Leakage}

\bookgraphic{0.453}{part-2-23-0001-011-04}{IS 得分异常跃升而 OOS 未同步改善，值得检查该次修改。}

IS

OOS-P

OOS-Q

四种假想曲线均用于提出诊断方向，不构成对具体退化机制的因果证明。特别是 Leakage 无法仅通过曲线形状确认。

这套实验设计在 Synthetic Environment 中相对容易实施，因为我们可以主动控制数据生成机制。

在真实金融环境中，严格构造来自相同分布的两个独立数据集则更加困难。

这也进一步说明了为什么 Synthetic Experiment 对该研究具有价值：它能够隔离现实数据中难以分离的统计机制。

\hypertarget{part-2-23-0001-diagnostic-toolbox}{%
\subsubsection{7.2 Diagnostic Toolbox}\label{part-2-23-0001-diagnostic-toolbox}}

\begin{readertable}
\begin{longtable}[]{@{}
  >{\raggedright\arraybackslash}p{(\columnwidth - 4\tabcolsep) * \real{0.34}}
  >{\raggedright\arraybackslash}p{(\columnwidth - 4\tabcolsep) * \real{0.30}}
  >{\raggedright\arraybackslash}p{(\columnwidth - 4\tabcolsep) * \real{0.32}}@{}}
\toprule
Method & 诊断对象 & 局限性 \\
\midrule
\endhead
Code \& Data Audit & Data Leakage & 需要完整代码和数据来源 \\
KS Test & 单变量分布变化 & 不能单独识别 Concept Drift \\
MMD / Two-sample Test & 多变量分布差异 & 结果受样本量等因素影响 \\
Domain Classifier & 区分源环境与目标环境 & 不能单独确定偏移机制 \\
Independent Holdout & Statistical Overfitting & 需要足够的独立数据 \\
Controlled Regime Simulation & Distribution Shift & 外部有效性有限 \\
Rollback / Ablation & 特定 Revision 的影响 & 需要控制其他修改及随机性 \\
\bottomrule
\end{longtable}
\end{readertable}

其中 KS Test（Kolmogorov--Smirnov Test）与你之前学习的 CDF 有直接联系。

其双样本检验统计量为：
\[D_{n,m}
=
\sup_x
|\widehat F_n(x)-\widehat F_m(x)|.\]

它衡量两个经验分布函数之间的最大差异。

但对于具有时间依赖性的金融数据，普通 IID KS Test 的显著性计算不能直接照搬。

\hypertarget{part-2-23-0001-revision-classification-ux7684ux8fdbux4e00ux6b65ux5b8cux5584}{%
\subsection{8. Revision Classification 的进一步完善}\label{part-2-23-0001-revision-classification-ux7684ux8fdbux4e00ux6b65ux5b8cux5584}}

结合 II-C、II-D、II-E 的知识，我们可以从统计机制角度扩展 Part I 中的 Revision Taxonomy。

\begin{readertable}
\begin{longtable}[]{@{}
  >{\raggedright\arraybackslash}p{(\columnwidth - 2\tabcolsep) * \real{0.48}}
  >{\raggedright\arraybackslash}p{(\columnwidth - 2\tabcolsep) * \real{0.48}}@{}}
\toprule
Revision Category & Potential Failure Mechanism \\
\midrule
\endhead
Feature Expansion & Classical Overfitting \\
Candidate Search Expansion & Selection Bias \\
Feedback and Memory Expansion & Adaptive Overfitting \\
Future-dependent Features & Look-Ahead Bias \\
Global Data Preprocessing & Preprocessing Leakage \\
Regime-specific Rules & Distribution Shift Sensitivity \\
Evaluation Metric Modification & Metric Gaming \\
Regularization and Robustness & 可能缓解性能退化 \\
\bottomrule
\end{longtable}
\end{readertable}

未来分析 Agent 的代码版本差异时，不应当仅仅记录：

\begin{quote}
Agent 增加了一个 Feature。
\end{quote}

还应该进一步记录这个 Feature 是否具有时间可用性、是否依赖历史评估反馈、是否针对特定市场环境，以及它是否扩大了搜索空间。

这会使我们的 Revision Classification 从单纯的软件功能分类，发展为具有统计学解释能力的分类体系。

\hypertarget{part-2-23-0001-literature-index}{%
\subsection{9. Literature Index}\label{part-2-23-0001-literature-index}}

本阶段不需要大量阅读论文。以下三篇已核对正式出版记录，均属于同行评审研究。

\begin{readertable}
\begin{longtable}[]{@{}
  >{\raggedright\arraybackslash}p{(\columnwidth - 4\tabcolsep) * \real{0.18}}
  >{\raggedright\arraybackslash}p{(\columnwidth - 4\tabcolsep) * \real{0.39}}
  >{\raggedright\arraybackslash}p{(\columnwidth - 4\tabcolsep) * \real{0.40}}@{}}
\toprule
Priority & Paper & Research Connection \\
\midrule
\endhead
P0 & Kaufman et al. (2012) & Data Leakage \\
P0 & Moreno-Torres et al. (2012) & Distribution Shift Taxonomy \\
P2 & Lipton et al. (2018) & Label Shift Correction \\
\bottomrule
\end{longtable}
\end{readertable}

\textbf{Kaufman, S., Rosset, S., Perlich, C. \& Stitelman, O. (2012).} Leakage in Data Mining: Formulation, Detection, and Avoidance. \emph{ACM Transactions on Knowledge Discovery from Data}. \href{https://doi.org/10.1145/2382577.2382579}{DOI: 10.1145/2382577.2382579}。

\href{https://doi.org/10.1145/2382577.2382579?trk=article-ssr-frontend-pulse_little-text-block}{来源 1} \href{https://cris.tau.ac.il/en/publications/leakage-in-data-mining-formulation-detection-and-avoidance/}{来源 2}

\textbf{Moreno-Torres, J. G., Raeder, T., Alaiz-Rodríguez, R., Chawla, N. V. \& Herrera, F. (2012).} A Unifying View on Dataset Shift in Classification. \emph{Pattern Recognition}, 45(1), 521--530. \href{https://www.sciencedirect.com/science/article/pii/S0031320311002901}{正式出版页面}。

\href{https://www.sciencedirect.com/science/article/pii/S0031320311002901}{来源 1}

\textbf{Lipton, Z. C., Wang, Y.-X. \& Smola, A. J. (2018).} Detecting and Correcting for Label Shift with Black Box Predictors. \emph{ICML 2018}. \href{https://proceedings.mlr.press/v80/lipton18a.html}{PMLR 正式会议论文}。

\href{https://proceedings.mlr.press/v80/lipton18a}{来源 1}

优先掌握前两篇的概念分类即可，暂时不需要深入第三篇的具体算法。

\hypertarget{part-2-23-0001-ii-e-conclusion-index}{%
\subsection{10. II-E --- Conclusion Index}\label{part-2-23-0001-ii-e-conclusion-index}}

\hypertarget{part-2-23-0001-ux6838ux5fc3ux7ed3ux8bba}{%
\subsubsection{10.1 核心结论}\label{part-2-23-0001-ux6838ux5fc3ux7ed3ux8bba}}

\begin{readertable}
\begin{longtable}[]{@{}
  >{\raggedright\arraybackslash}p{(\columnwidth - 4\tabcolsep) * \real{0.12}}
  >{\raggedright\arraybackslash}p{(\columnwidth - 4\tabcolsep) * \real{0.71}}
  >{\raggedright\arraybackslash}p{(\columnwidth - 4\tabcolsep) * \real{0.13}}@{}}
\toprule
ID & Key Result & Index \\
\midrule
\endhead
E1 & Data Leakage 的本质是违反信息可用性或实验协议 & §1 \\
E2 & 金融预测必须基于预测时刻已经可获得的信息 & §1.1 \\
E3 & Data Leakage 不等于 Adaptive Overfitting & §2 \\
E4 & Covariate Shift 改变输入分布 & §3.2 \\
E5 & Label Shift 改变标签边际分布 & §3.3 \\
E6 & Concept Drift 改变输入与输出之间的统计关系 & §3.4 \\
E7 & 即使完全没有过拟合，Distribution Shift 也可以造成性能损失 & §4 \\
E8 & 观察到的 Gap 包含统计误差、分布变化及 OOS 评估噪声 & §5 \\
E9 & Across Iterations 与 Across Environments 是不同研究维度 & §6 \\
E10 & 同分布与跨分布对照实验有助于诊断退化机制 & §7 \\
\bottomrule
\end{longtable}
\end{readertable}

\hypertarget{part-2-23-0001-core-formula-sheet}{%
\subsubsection{10.2 Core Formula Sheet}\label{part-2-23-0001-core-formula-sheet}}

\textbf{Result 1 --- Information Availability Constraint}

\[\hat r_{t+1}\text{ is }\mathcal F_t\text{-measurable}.\]

\textbf{Result 2 --- Covariate Shift Correction}

\[R_Q(h)
=
\mathbb E_P
\left[
\frac{q(X)}{p(X)}
\ell(h(X),Y)
\right].\]

\textbf{Result 3 --- Concept Drift Regression Example}

\[R_Q(h)-R_P(h)=\Delta^2.\]

\textbf{Result 4 --- Observed Generalization Gap Decomposition}

\[\boxed{
\begin{aligned}
\widehat R_m^Q(h)-\widehat R_n(h)
={}&G_{\mathrm{stat}}(h)\\
&+G_{\mathrm{shift}}(h)\\
&+\eta_{\mathrm{OOS}}(h),
\end{aligned}
}\]

其中：
\[\eta_{\mathrm{OOS}}(h)
=
\widehat R_m^Q(h)-R_Q(h).\]

以上公式均应结合对应章节的成立条件理解。

\hypertarget{part-2-23-0001-part-ii-consolidated-diagnostic-framework}{%
\subsection{11. Part II --- Consolidated Diagnostic Framework}\label{part-2-23-0001-part-ii-consolidated-diagnostic-framework}}

现在，我们已经积累了五个模块的统计理论，可以将其整合为一张真正具有科研诊断意义的知识图谱。

\textbf{Observed OOS Performance Decay}

需要解释的实验现象

↓

\textbf{Classical Overfitting}

Model Complexity

II-A / II-B

\textbf{Selection Bias}

Multiple Comparisons

II-C

\textbf{Adaptive Overfitting}

Feedback Information

II-D

\textbf{Data Leakage}

Invalid Information

II-E

\textbf{Distribution Shift}

Changing Environment

II-E

\textbf{Evaluation Noise}

Finite-sample Uncertainty

II-B / II-C

↓

\textbf{Controlled Experiments \& Ablation Studies}

逐步隔离不同机制，建立可信的诊断结论

这里的分类体系及实验方案是我们根据项目描述建立的候选研究框架，尚未通过实际实验验证，也不代表课题组已经采用这些方法。

\textbf{至此，我们基本掌握了该项目所需要的核心统计机制。}

\hypertarget{part-2-23-0001-next-ii-f-time-series-cross-validation-and-backtesting}{%
\subsection{Next --- II-F: Time-Series Cross-Validation and Backtesting}\label{part-2-23-0001-next-ii-f-time-series-cross-validation-and-backtesting}}

II-A --- Risk and Generalization

II-B --- Concentration Inequalities

II-C --- Model Selection Bias \& Multiple Testing

II-D --- Adaptive Overfitting \& Holdout Reuse

II-E --- Data Leakage \& Distribution Shift

II-F --- Time-Series Validation \& Backtesting

下一阶段是 Part II 的最后一个模块，也是从 Statistical Learning Theory 过渡到实际金融实验的桥梁。

\begin{readertable}
\begin{longtable}[]{@{}
  >{\raggedright\arraybackslash}p{(\columnwidth - 2\tabcolsep) * \real{0.38}}
  >{\raggedright\arraybackslash}p{(\columnwidth - 2\tabcolsep) * \real{0.58}}@{}}
\toprule
Topic & 核心问题 \\
\midrule
\endhead
Ordinary K-fold CV & 为什么普通交叉验证可能不适合金融预测？ \\
Expanding Window & 如何利用不断积累的历史数据？ \\
Rolling Window & 如何适应非平稳环境？ \\
Walk-forward Validation & 如何模拟真实预测过程？ \\
Purged Cross-Validation & 如何避免重叠标签造成的信息泄漏？ \\
Embargo & 如何处理相邻时间窗口的信息依赖？ \\
Backtesting & 如何可靠评价历史交易策略？ \\
Nested Validation & 如何分离模型训练、Agent 优化与最终评价？ \\
\bottomrule
\end{longtable}
\end{readertable}

下一阶段将主要关注 \textbf{Experimental Protocol}，减少数学证明，重点理解不同数据划分方式的动机、成立条件和适用范围。

完成 II-F 后，我们会对整个 Part II 进行一次统一的 Conclusion Index 整理，再正式进入 Part III --- Agentic AI and Recursive Self-Improvement。
